\documentclass[11pt,a4paper]{article}
\usepackage[utf8]{inputenc}
\usepackage[T1]{fontenc}
\usepackage[french]{babel}
\usepackage{amsmath,amssymb,amsfonts,amsthm,mathtools}
\usepackage[margin=1.8cm,top=2.2cm,bottom=2.2cm,headheight=14pt]{geometry}
\usepackage{enumitem}
\usepackage{xcolor}
\usepackage{tcolorbox}
\usepackage{tikz}
\usepackage{fancyhdr}
\usepackage{lastpage}
\usepackage{tabularx}
\usepackage{array}
\usepackage{booktabs}
\usepackage{hyperref}

\tcbuselibrary{skins,breakable}
\usetikzlibrary{calc,angles,quotes,shapes.geometric,patterns,arrows.meta}

% ----------------------------------------------------
% PALETTE DE COULEURS ÉLÉGANTE & SOBRE (STYLE MPSI / LLG)
% ----------------------------------------------------
\definecolor{primary}{RGB}{18, 48, 100}       % Bleu marine royal
\definecolor{secondary}{RGB}{130, 25, 35}     % Grenat / Carmin
\definecolor{forestgreen}{RGB}{24, 110, 50}   % Vert émeraude sombre
\definecolor{darkamber}{RGB}{170, 95, 10}     % Ambre doré
\definecolor{royalpurple}{RGB}{85, 35, 135}   % Pourpre / Violet MPSI
\definecolor{lightgraybg}{RGB}{248, 250, 252} % Gris bleuté ultra clair
\definecolor{linegray}{RGB}{215, 222, 230}    % Filet de séparation

% ----------------------------------------------------
% ENVIRONNEMENTS TCOLORBOX
% ----------------------------------------------------
% Boîte Énoncé
\newtcolorbox{enoncebox}[1][]{
  enhanced,
  colback=lightgraybg,
  colframe=primary,
  coltitle=white,
  fonttitle=\bfseries\sffamily,
  title={#1},
  arc=1.5mm,
  boxrule=0.8pt,
  left=4mm, right=4mm, top=3mm, bottom=3mm,
  breakable
}

% Boîte Corrigé
\newtcolorbox{correctionbox}[1][]{
  enhanced,
  colback=forestgreen!4!white,
  colframe=forestgreen!85!black,
  coltitle=white,
  fonttitle=\bfseries\sffamily,
  title={#1},
  arc=1.5mm,
  boxrule=0.8pt,
  left=4mm, right=4mm, top=3mm, bottom=3mm,
  breakable
}

% Boîte Méthode & Vigilance (Conseils du correcteur LLG)
\newtcolorbox{vigilancebox}[1][]{
  enhanced,
  colback=darkamber!5!white,
  colframe=darkamber!90!black,
  coltitle=white,
  fonttitle=\bfseries\sffamily\small,
  title={\textsc{Rigueur \& Rédaction : Les Pièges Classiques} \hfill #1},
  arc=1.2mm,
  boxrule=0.7pt,
  left=3.5mm, right=3.5mm, top=2.5mm, bottom=2.5mm,
  breakable
}

% Boîte Prolongement Mathématique (Niveau Prépa MPSI / Concours Général)
\newtcolorbox{prolongementbox}[1][]{
  enhanced,
  colback=royalpurple!4!white,
  colframe=royalpurple!85!black,
  coltitle=white,
  fonttitle=\bfseries\sffamily,
  title={\textsc{Prolongement Mathématique \& Culturel (MPSI / Approfondissement)} \hfill #1},
  arc=1.5mm,
  boxrule=0.8pt,
  left=4mm, right=4mm, top=3mm, bottom=3mm,
  breakable
}

% Badge de barème
\newcommand{\bareme}[1]{\hfill\textbf{\textcolor{secondary}{\sffamily[#1]}}}
\newcommand{\pts}[1]{\textbf{\textcolor{secondary}{\sffamily(#1)}}}

% Macros mathématiques usuelles
\newcommand{\N}{\mathbb{N}}
\newcommand{\Z}{\mathbb{Z}}
\newcommand{\Q}{\mathbb{Q}}
\newcommand{\R}{\mathbb{R}}
\newcommand{\C}{\mathbb{C}}
\newcommand{\dsp}{\displaystyle}

% ----------------------------------------------------
% CONFIGURATION DES EN-TÊTES ET PIEDS DE PAGE
% ----------------------------------------------------
\pagestyle{fancy}
\fancyhf{}
\lhead{\footnotesize\sffamily\textbf{Lycée Louis-le-Grand} --- Classe de 1\iere{}E3 (Spécialité Maths)}
\chead{\footnotesize\sffamily\textbf{DST Récurrence} --- Énoncé \& Corrigé Détaillé}
\rhead{\footnotesize\sffamily 16 décembre 2021}
\lfoot{\footnotesize\sffamily Soutien Scolaire d'Excellence --- Mathématiques}
\rfoot{\footnotesize\sffamily Page \thepage{} / \pageref{LastPage}}
\renewcommand{\headrulewidth}{0.5pt}
\renewcommand{\footrulewidth}{0.4pt}

\hypersetup{
  colorlinks=true,
  linkcolor=primary,
  urlcolor=primary,
  citecolor=secondary,
  pdfauthor={Enseignant MPSI / Soutien Scolaire LLG},
  pdftitle={DST Récurrence 1E3 - Sujet et Corrigé Intégral}
}

% ----------------------------------------------------
% DÉBUT DU DOCUMENT
% ----------------------------------------------------
\begin{document}

% ----------------------------------------------------
% CARTOUCHE DE PRÉSENTATION OFFICIEL
% ----------------------------------------------------
\begin{center}
  {\fontsize{17}{21}\selectfont\textbf{\textsc{Devoir Sur Table n°3 : Le Raisonnement par Récurrence}}}\\[4pt]
  {\large\sffamily\textbf{Classe de Première 3 (Spécialité Mathématiques) --- 16 décembre 2021}}\\[2pt]
  {\small\sffamily Niveau : Lycée Louis-le-Grand / Préparation aux Classes Préparatoires (MPSI-PCSI)}\\[6pt]
  \rule{0.9\textwidth}{0.8pt}
\end{center}

\vspace{-0.3cm}

\begin{tcolorbox}[enhanced, colback=white, colframe=linegray, boxrule=0.6pt, arc=1mm]
\begin{tabularx}{\textwidth}{lXr}
  \textbf{Durée de l'épreuve :} 2 heures & \textbf{Calculatrice :} Interdite (calcul formel et arithmétique manuels) & \textbf{Total :} 20 points + 2 bonus \\
  \textbf{Notions testées :} & \multicolumn{2}{l}{Suites récurrentes, inégalités, géométrie pythagoricienne, sommes télescopiques, arithmétique.}
\end{tabularx}
\end{tcolorbox}

\vspace{0.1cm}

% Tableau récapitulatif du barème
\begin{center}
\renewcommand{\arraystretch}{1.2}
\small
\begin{tabularx}{\textwidth}{|c|X|c|c|}
\hline
\textbf{Exercice} & \textbf{Objet d'étude \& Thématique mathématique} & \textbf{Questions} & \textbf{Barème} \\
\hline
\textbf{Exercice 1} & Suite arithmético-géométrique : monotonie, minoration et terme général & 3 & \textbf{4 points} \\
\hline
\textbf{Exercice 2} & Géométrie euclidienne : la Spirale de Théodore de Cyrène ($\sqrt{n}$) & 3 & \textbf{4 points} \\
\hline
\textbf{Exercice 3} & Sommation cubique, conjecture inductive et formule du binôme & 3 & \textbf{4 points} \\
\hline
\textbf{Exercice 4} & Somme alternée d'entiers $\sum_{i=1}^n (-1)^i i$ et factorisations fines & 1 & \textbf{4 points} \\
\hline
\textbf{Exercice 5} & Arithmétique décimale : puissances de $10$ modulo $3$ et somme des chiffres & 2 & \textbf{4 points} \\
\hline
\textbf{Bonus} & Démonstration directe non inductive par télescopage cubique & 1 & \textbf{+2 points} \\
\hline
\multicolumn{3}{|r|}{\textbf{Total général indicatif}} & \textbf{20 pts (+ 2 bonus)} \\
\hline
\end{tabularx}
\end{center}

\vspace{0.3cm}

\begin{center}
  \large\textbf{\textsc{Partie I : Énoncés Complets du Devoir}}
\end{center}

% ----------------------------------------------------
% ÉNONCÉ EXERCICE 1
% ----------------------------------------------------
\begin{enoncebox}[Exercice 1 \bareme{4 points}]
Soit $(u_n)$ la suite numérique définie par :
\[
\begin{cases}
u_0 = 6 \\
u_{n+1} = \dfrac{1}{2}u_n - 1 \quad (n \in \N)
\end{cases}
\]
\begin{enumerate}[label=\textbf{\arabic*.}]
    \item Montrer que la suite $(u_n)$ est décroissante. \bareme{1,25 pt}
    \item Montrer que la suite est minorée par $-2$. \bareme{1,25 pt}
    \item Montrer que pour tout $n \in \N$, $u_n = 2\left(2^{2-n} - 1\right)$. \bareme{1,5 pt}
\end{enumerate}
\end{enoncebox}

% ----------------------------------------------------
% ÉNONCÉ EXERCICE 2
% ----------------------------------------------------
\begin{enoncebox}[Exercice 2 --- La Spirale de Théodore de Cyrène \bareme{4 points}]
Dans la figure ci-dessous, le triangle $OA_1A_2$ est un triangle rectangle isocèle en $A_1$ tel que $OA_1 = 1$.\\
Pour $n \in \N^*$, on construit des triangles rectangles $OA_n A_{n+1}$ successifs d'hypothénuses $OA_{n+1}$ avec $A_k A_{k+1} = 1$ comme sur la figure ci-dessous.

\begin{center}
\begin{tikzpicture}[scale=1.2, >=latex, line cap=round, line join=round]
  % Point d'origine O
  \coordinate (O) at (0,0);
  \fill[black] (O) circle (1.8pt);
  \node[below=3pt, font=\small] at (O) {$O$};

  % Coordonnées exactes des sommets A1 à A7
  \coordinate (A1) at (1.800, 0.000);
  \coordinate (A2) at (1.800, 1.800);
  \coordinate (A3) at (0.527, 3.073);
  \coordinate (A4) at (-1.247, 3.377);
  \coordinate (A5) at (-2.935, 2.754);
  \coordinate (A6) at (-4.167, 1.441);
  \coordinate (A7) at (-4.755, -0.260);

  % Remplissages discrets des triangles
  \fill[primary!7] (O) -- (A1) -- (A2) -- cycle;
  \fill[primary!12] (O) -- (A2) -- (A3) -- cycle;
  \fill[primary!7] (O) -- (A3) -- (A4) -- cycle;
  \fill[primary!12] (O) -- (A4) -- (A5) -- cycle;
  \fill[primary!7] (O) -- (A5) -- (A6) -- cycle;

  % Rayons / Hypoténuses OA_k
  \draw[thick, primary!80!black] (O) -- (A1) node[midway, below=2pt, font=\footnotesize] {$1$};
  \draw[thick, primary!80!black] (O) -- (A2) node[midway, above left=-1pt, font=\scriptsize] {$OA_2$};
  \draw[thick, primary!80!black] (O) -- (A3);
  \draw[thick, primary!80!black] (O) -- (A4);
  \draw[thick, primary!80!black] (O) -- (A5);
  \draw[thick, primary!80!black] (O) -- (A6);
  \draw[dashed, thick, gray!60] (O) -- (A7);

  % Côtés extérieurs orthogonaux de longueur 1
  \draw[very thick, secondary] (A1) -- (A2) node[midway, right=2pt, font=\footnotesize] {$1$};
  \draw[very thick, secondary] (A2) -- (A3) node[midway, above right=1pt, font=\footnotesize] {$1$};
  \draw[very thick, secondary] (A3) -- (A4) node[midway, above=2pt, font=\footnotesize] {$1$};
  \draw[very thick, secondary] (A4) -- (A5) node[midway, above left=1pt, font=\footnotesize] {$1$};
  \draw[very thick, secondary] (A5) -- (A6) node[midway, left=2pt, font=\footnotesize] {$1$};
  \draw[dashed, very thick, secondary!60] (A6) -- (A7);

  % Symboles d'angle droit en A1, A2, A3, A4, A5
  \draw[black!60, thin] (1.55, 0) -- (1.55, 0.25) -- (1.8, 0.25);
  \draw[black!60, thin] ($(A2)!3.2mm!(O)$) -- ($($(A2)!3.2mm!(O)$)+($(A2)!3.2mm!(A3)$)-(A2)$) -- ($(A2)!3.2mm!(A3)$);
  \draw[black!60, thin] ($(A3)!3.2mm!(O)$) -- ($($(A3)!3.2mm!(O)$)+($(A3)!3.2mm!(A4)$)-(A3)$) -- ($(A3)!3.2mm!(A4)$);
  \draw[black!60, thin] ($(A4)!3.2mm!(O)$) -- ($($(A4)!3.2mm!(O)$)+($(A4)!3.2mm!(A5)$)-(A4)$) -- ($(A4)!3.2mm!(A5)$);
  \draw[black!60, thin] ($(A5)!3.2mm!(O)$) -- ($($(A5)!3.2mm!(O)$)+($(A5)!3.2mm!(A6)$)-(A5)$) -- ($(A5)!3.2mm!(A6)$);

  % Sommets
  \fill[black] (A1) circle (1.6pt) node[below right, font=\small] {$A_1$};
  \fill[black] (A2) circle (1.6pt) node[right=2pt, font=\small] {$A_2$};
  \fill[black] (A3) circle (1.6pt) node[above right=1pt, font=\small] {$A_3$};
  \fill[black] (A4) circle (1.6pt) node[above=2pt, font=\small] {$A_4$};
  \fill[black] (A5) circle (1.6pt) node[above left=1pt, font=\small] {$A_5$};
  \fill[black] (A6) circle (1.6pt) node[left=2pt, font=\small] {$A_6$};

  \node[font=\large] at (-3.2, 0.2) {$\dots$};
\end{tikzpicture}
\end{center}

\noindent Conjecturer puis démontrer l'expression de $OA_n$ en fonction de $n$ pour $n \in \N^*$.
\end{enoncebox}

% ----------------------------------------------------
% ÉNONCÉ EXERCICE 3
% ----------------------------------------------------
\begin{enoncebox}[Exercice 3 \bareme{4 points}]
Soit $(u_n)$ la suite définie par :
\[
\begin{cases}
u_0 = 0 \\
u_{n+1} = u_n + 3n(n+1) + 1 \quad (n \in \N)
\end{cases}
\]
\begin{enumerate}[label=\textbf{\arabic*.}]
    \item Calculer $u_1$, $u_2$ et $u_3$. \bareme{1 pt}
    \item Conjecturer l'expression de $u_n$ en fonction de $n$. \bareme{0,5 pt}
    \item Démontrer la conjecture par récurrence. \bareme{2,5 pts}
\end{enumerate}
\end{enoncebox}

% ----------------------------------------------------
% ÉNONCÉ EXERCICE 4
% ----------------------------------------------------
\begin{enoncebox}[Exercice 4 \bareme{4 points}]
Montrer que pour tout $n \in \N^*$ :
\[
\sum_{i=1}^n (-1)^i \, i = \frac{(-1)^n (2n+1) - 1}{4}
\]
\end{enoncebox}

% ----------------------------------------------------
% ÉNONCÉ EXERCICE 5
% ----------------------------------------------------
\begin{enoncebox}[Exercice 5 \bareme{4 points}]
\begin{enumerate}[label=\textbf{\arabic*.}]
    \item Montrer que pour tout entier naturel $n$, il existe un entier naturel $b_n$ tel que :
    \[
    10^n = 3b_n + 1
    \] \bareme{2 pts}
    \item En déduire que si $a_0, a_1, \dots, a_p$ sont des entiers naturels, alors :
    \[
    \left( 3 \text{ divise } \sum_{k=0}^p a_k 10^k \right) \iff \left( 3 \text{ divise } \sum_{k=0}^p a_k \right)
    \]
    Interprétez cette équivalence arithmétique. \bareme{2 pts}
\end{enumerate}
\end{enoncebox}

% ----------------------------------------------------
% ÉNONCÉ DU BONUS
% ----------------------------------------------------
\begin{enoncebox}[Bonus (+2 points)]
Démontrer la conjecture de l'exercice 3 de façon directe (sans raisonnement par récurrence). \bareme{+2 pts}
\end{enoncebox}

\newpage

% ----------------------------------------------------
% PARTIE II : CORRIGÉS DÉTAILLÉS & MÉTHODOLOGIE
% ----------------------------------------------------
\begin{center}
  \large\textbf{\textsc{Partie II : Corrigé Exhaustif, Barème Détaillé \& Prolongements}}
\end{center}

\vspace{0.2cm}

% ----------------------------------------------------
% CORRIGÉ EXERCICE 1
% ----------------------------------------------------
\subsection*{Correction détaillée de l'Exercice 1}

\begin{correctionbox}[Corrigé de l'Exercice 1 \bareme{Total : 4 points}]
La suite $(u_n)_{n \in \N}$ est définie par $u_0 = 6$ et pour tout $n \in \N$, $u_{n+1} = \dfrac{1}{2}u_n - 1$.

\begin{enumerate}[label=\textbf{\arabic*.}]
    \item \textbf{Démontrons par récurrence que la suite $(u_n)$ est strictement décroissante.} \bareme{1,25 pt}

    Pour tout $n \in \N$, introduisons le prédicat :
    \[
    \mathcal{P}(n) : \quad \ll u_n > u_{n+1} \gg
    \]
    
    \begin{itemize}[leftmargin=*, label=$\bullet$]
        \item \textbf{Initialisation ($n = 0$) :} \pts{0,25 pt}\\
        Le premier terme est $u_0 = 6$.\\
        Calculons le terme suivant grâce à la relation de récurrence :
        \[
        u_1 = \frac{1}{2} u_0 - 1 = \frac{1}{2} \times 6 - 1 = 3 - 1 = 2
        \]
        Puisque $6 > 2$, l'inégalité $u_0 > u_1$ est satisfaite.\\
        Ainsi, la propriété $\mathcal{P}(0)$ est vraie.

        \item \textbf{Hérédité :} \pts{0,75 pt}\\
        Soit $k \in \N$ un entier fixé quelconque.\\
        Supposons que la propriété $\mathcal{P}(k)$ soit vraie, c'est-à-dire que $u_k > u_{k+1}$ (\textbf{Hypothèse de Récurrence}).\\
        Montrons que la propriété $\mathcal{P}(k+1)$ est vraie, c'est-à-dire que $u_{k+1} > u_{k+2}$.\\
        
        Partons de l'hypothèse de récurrence :
        \[
        u_k > u_{k+1}
        \]
        La fonction affine $x \mapsto \dfrac{1}{2}x - 1$ a pour coefficient directeur $\dfrac{1}{2} > 0$ ; elle est donc strictement croissante sur $\R$. En multipliant les deux membres par le réel strictement positif $\dfrac{1}{2}$ :
        \[
        \frac{1}{2}u_k > \frac{1}{2}u_{k+1}
        \]
        En retranchant le réel $1$ à chaque membre de l'inégalité :
        \[
        \frac{1}{2}u_k - 1 > \frac{1}{2}u_{k+1} - 1
        \]
        Or, par définition même de la suite récurrente :
        \[
        \frac{1}{2}u_k - 1 = u_{k+1} \quad \text{et} \quad \frac{1}{2}u_{k+1} - 1 = u_{k+2}
        \]
        On en déduit donc immédiatement :
        \[
        u_{k+1} > u_{k+2}
        \]
        Le prédicat $\mathcal{P}(k+1)$ est ainsi démontré. L'hérédité est établie pour tout $k \in \N$.

        \item \textbf{Conclusion :} \pts{0,25 pt}\\
        La propriété $\mathcal{P}(0)$ est initialisée et pour tout $k \in \N$, $\mathcal{P}(k) \implies \mathcal{P}(k+1)$.\\
        D'après le principe de récurrence simple, $\mathcal{P}(n)$ est vraie pour tout tout entier naturel $n \in \N$.\\
        Par conséquent, $\forall n \in \N, \; u_n > u_{n+1}$ : la suite $(u_n)_{n \in \N}$ est \textbf{strictement décroissante}.
    \end{itemize}

    \item \textbf{Démontrons par récurrence que la suite $(u_n)$ est minorée par $-2$.} \bareme{1,25 pt}

    Pour tout $n \in \N$, définissons le prédicat :
    \[
    \mathcal{Q}(n) : \quad \ll u_n > -2 \gg
    \]

    \begin{itemize}[leftmargin=*, label=$\bullet$]
        \item \textbf{Initialisation ($n = 0$) :} \pts{0,25 pt}\\
        $u_0 = 6$. Comme $6 > -2$, le prédicat $\mathcal{Q}(0)$ est rigoureusement vérifié.

        \item \textbf{Hérédité :} \pts{0,75 pt}\\
        Soit $k \in \N$. Supposons que $\mathcal{Q}(k)$ soit vraie, c'est-à-dire que $u_k > -2$ (\textbf{H.R.}).\\
        Montrons que $\mathcal{Q}(k+1)$ est vraie, à savoir $u_{k+1} > -2$.\\
        
        En multipliant l'inégalité de l'hypothèse de récurrence par $\dfrac{1}{2} > 0$ :
        \[
        u_k > -2 \implies \frac{1}{2}u_k > \frac{1}{2} \times (-2) \iff \frac{1}{2}u_k > -1
        \]
        En soustrayant $1$ aux deux membres :
        \[
        \frac{1}{2}u_k - 1 > -1 - 1 \iff \frac{1}{2}u_k - 1 > -2
        \]
        Or, le membre de gauche correspond exactement à $u_{k+1}$. Il s'ensuit que :
        \[
        u_{k+1} > -2
        \]
        L'implication $\mathcal{Q}(k) \implies \mathcal{Q}(k+1)$ est donc vérifiée pour tout $k \in \N$.

        \item \textbf{Conclusion :} \pts{0,25 pt}\\
        Le principe de récurrence garantit que pour tout entier naturel $n \in \N$, $u_n > -2$.\\
        La suite $(u_n)_{n \in \N}$ est donc \textbf{strictement minorée par $-2$}.
    \end{itemize}

    \item \textbf{Démontrons par récurrence que pour tout $n \in \N$, $u_n = 2\left(2^{2-n} - 1\right)$.} \bareme{1,5 pt}

    Pour tout $n \in \N$, notons le prédicat :
    \[
    \mathcal{R}(n) : \quad \ll u_n = 2\left(2^{2-n} - 1\right) \gg
    \]
    Remarquons que cette formule s'écrit de manière équivalente :
    \[
    u_n = 2 \cdot 2^{2-n} - 2 = 2^{3-n} - 2 = 8 \left(\frac{1}{2}\right)^n - 2
    \]

    \begin{itemize}[leftmargin=*, label=$\bullet$]
        \item \textbf{Initialisation ($n = 0$) :} \pts{0,25 pt}\\
        D'une part, par donnée de l'énoncé, le membre de gauche vaut : $u_0 = 6$.\\
        D'autre part, en évaluant l'expression du membre de droite pour $n = 0$ :
        \[
        2\left(2^{2-0} - 1\right) = 2\left(2^2 - 1\right) = 2(4 - 1) = 2 \times 3 = 6
        \]
        Les deux membres coïncident parfaitement : $u_0 = 2\left(2^{2-0} - 1\right)$.\\
        L'initialisation $\mathcal{R}(0)$ est donc avérée.

        \item \textbf{Hérédité :} \pts{1,0 pt}\\
        Soit $k \in \N$. Supposons que $\mathcal{R}(k)$ soit vraie, c'est-à-dire que :
        \[
        u_k = 2\left(2^{2-k} - 1\right) \quad \text{(\textbf{H.R.})}
        \]
        Montrons que $\mathcal{R}(k+1)$ est vraie, c'est-à-dire que :
        \[
        u_{k+1} = 2\left(2^{2-(k+1)} - 1\right) = 2\left(2^{1-k} - 1\right)
        \]
        Injectons l'hypothèse de récurrence dans la relation définissant $u_{k+1}$ :
        \begin{align*}
        u_{k+1} &= \frac{1}{2} u_k - 1 \\
                &= \frac{1}{2} \left[ 2\left(2^{2-k} - 1\right) \right] - 1 \\
                &= \left(2^{2-k} - 1\right) - 1 \quad \text{(car $\frac{1}{2} \times 2 = 1$)} \\
                &= 2^{2-k} - 2
        \end{align*}
        Mettons à présent le facteur $2$ en évidence :
        \[
        u_{k+1} = 2\left(2^{1-k} - 1\right) = 2\left(2^{2-(k+1)} - 1\right)
        \]
        Ce qui correspond très exactement à la proposition $\mathcal{R}(k+1)$. L'hérédité est prouvée.

        \item \textbf{Conclusion :} \pts{0,25 pt}\\
        Par le principe de récurrence, pour tout $n \in \N$ :
        \[
        u_n = 2\left(2^{2-n} - 1\right) = 2^{3-n} - 2 = 8 \left(\frac{1}{2}\right)^n - 2
        \]
    \end{itemize}
\end{enumerate}
\end{correctionbox}

\begin{vigilancebox}
\textbf{Erreur classique n°1 :} Ne jamais écrire en amorce d'hérédité « \emph{Supposons que pour tout $k$, $\mathcal{P}(k)$ est vraie} ». Cela reviendrait à supposer ce que l'on cherche à démontrer ! Il faut fixer \textbf{un} entier quelconque $k$ et supposer $\mathcal{P}(k)$ pour \textbf{cet} entier.\\
\textbf{Erreur classique n°2 :} Ne pas oublier de citer le sens de variation de la fonction affine sous-jacente ou de justifier que multiplier par $\dfrac{1}{2} > 0$ conserve le sens des inégalités.
\end{vigilancebox}

\begin{prolongementbox}[Théorie générale des suites arithmético-géométriques]
Toute suite de la forme $u_{n+1} = a u_n + b$ avec $a \neq 1$ admet un unique \textbf{point fixe} $\ell$, solution de l'équation $\ell = a \ell + b \iff \ell = \dfrac{b}{1-a}$.
Ici, avec $a = \frac{1}{2}$ et $b = -1$ :
\[
\ell = \frac{1}{2}\ell - 1 \iff \frac{1}{2}\ell = -1 \iff \ell = -2
\]
En posant la suite auxiliaire $v_n = u_n - \ell = u_n - (-2) = u_n + 2$, on a :
\[
v_{n+1} = u_{n+1} + 2 = \left(\frac{1}{2}u_n - 1\right) + 2 = \frac{1}{2}u_n + 1 = \frac{1}{2}(u_n + 2) = \frac{1}{2}v_n
\]
$(v_n)$ est donc une suite géométrique de raison $q = \frac{1}{2}$ et de premier terme $v_0 = u_0 + 2 = 6 + 2 = 8$.
Ainsi :
\[
v_n = 8 \left(\frac{1}{2}\right)^n \implies u_n = v_n - 2 = 8 \left(\frac{1}{2}\right)^n - 2 = 2\left(2^{2-n} - 1\right)
\]
Puisque $|q| = \frac{1}{2} < 1$, on a $\lim_{n \to +\infty} \left(\frac{1}{2}\right)^n = 0$, donc $\lim_{n \to +\infty} u_n = -2$.
Ce résultat confirme le \textbf{théorème de convergence monotone} : toute suite décroissante et minorée converge vers sa borne inférieure.
\end{prolongementbox}

\vspace{0.3cm}

% ----------------------------------------------------
% CORRIGÉ EXERCICE 2
% ----------------------------------------------------
\subsection*{Correction détaillée de l'Exercice 2}

\begin{correctionbox}[Corrigé de l'Exercice 2 \bareme{Total : 4 points}]
\textbf{Contexte géométrique :} Le premier triangle $OA_1A_2$ est rectangle en $A_1$ avec $OA_1 = 1$ et $A_1A_2 = 1$.\\
Pour tout $n \in \N^*$, le triangle $OA_n A_{n+1}$ est rectangle en $A_n$, d'hypoténuse $OA_{n+1}$, et son côté adjacent vérifie $A_n A_{n+1} = 1$.

\begin{enumerate}[label=\textbf{\arabic*.}]
    \item \textbf{Calcul des premières longueurs (premiers termes) :} \bareme{1,0 pt}
    \begin{itemize}[leftmargin=*, label=$\bullet$]
        \item Pour $n = 1$ : $OA_1 = 1 = \sqrt{1}$.
        \item Dans le triangle $OA_1A_2$ rectangle en $A_1$, le théorème de Pythagore s'énonce :
        \[
        OA_2^2 = OA_1^2 + A_1A_2^2 = 1^2 + 1^2 = 1 + 1 = 2
        \]
        Puisque $OA_2$ est une distance géométrique ($OA_2 > 0$) :
        \[
        OA_2 = \sqrt{2} \bareme{0,5 pt}
        \]
        \item Dans le triangle $OA_2A_3$ rectangle en $A_2$, le théorème de Pythagore fournit :
        \[
        OA_3^2 = OA_2^2 + A_2A_3^2 = \left(\sqrt{2}\right)^2 + 1^2 = 2 + 1 = 3
        \]
        Puisque $OA_3 > 0$ :
        \[
        OA_3 = \sqrt{3} \bareme{0,5 pt}
        \]
    \end{itemize}

    \item \textbf{Formulation de la conjecture :} \bareme{0,5 pt}\\
    En calculant le terme suivant, on trouve $OA_4 = \sqrt{OA_3^2 + 1^2} = \sqrt{3 + 1} = \sqrt{4} = 2$.\\
    Les premiers termes de la suite de longueurs $(OA_n)_{n \in \N^*}$ sont donc $\sqrt{1}, \sqrt{2}, \sqrt{3}, \sqrt{4}, \dots$\\
    On conjecture que pour tout entier $n \in \N^*$ :
    \[
    OA_n = \sqrt{n}
    \]

    \item \textbf{Démonstration de la conjecture par récurrence :} \bareme{1,5 pt}

    Pour tout $n \in \N^*$, définissons la propriété :
    \[
    \mathcal{P}(n) : \quad \ll OA_n = \sqrt{n} \gg
    \]

    \begin{itemize}[leftmargin=*, label=$\bullet$]
        \item \textbf{Initialisation ($n = 1$) :} \pts{0,25 pt}\\
        Pour $n = 1$, $OA_1 = 1$. Or $\sqrt{1} = 1$.\\
        L'égalité $OA_1 = \sqrt{1}$ est vérifiée : la propriété $\mathcal{P}(1)$ est vraie.

        \item \textbf{Hérédité :} \pts{1,0 pt}\\
        Soit $k \in \N^*$ un entier naturel fixé non nul.\\
        Supposons que $\mathcal{P}(k)$ soit vraie, à savoir $OA_k = \sqrt{k}$ (\textbf{Hypothèse de Récurrence}).\\
        Montrons que $\mathcal{P}(k+1)$ est vraie, c'est-à-dire que $OA_{k+1} = \sqrt{k+1}$.\\
        
        Considérons le triangle $OA_k A_{k+1}$. Par hypothèse géométrique de construction de la spirale, ce triangle est rectangle en $A_k$ et le côté $A_k A_{k+1}$ a pour longueur $1$.\\
        D'après le \textbf{théorème de Pythagore} appliqué à ce triangle :
        \[
        OA_{k+1}^2 = OA_k^2 + A_k A_{k+1}^2
        \]
        Or, d'après l'hypothèse de récurrence, $OA_k = \sqrt{k}$, d'où $OA_k^2 = \left(\sqrt{k}\right)^2 = k$.\\
        De plus, $A_k A_{k+1}^2 = 1^2 = 1$.\\
        En remplaçant ces valeurs dans l'égalité pythagoricienne :
        \[
        OA_{k+1}^2 = k + 1
        \]
        $OA_{k+1}$ désignant une longueur géométrique, c'est un nombre réel strictement positif. Par passage à la racine carrée :
        \[
        OA_{k+1} = \sqrt{k+1}
        \]
        La proposition $\mathcal{P}(k+1)$ est donc exacte. L'hérédité est validée pour tout $k \in \N^*$.

        \item \textbf{Conclusion :} \pts{0,25 pt}\\
        La propriété est vraie au rang initial $n = 1$ et est héréditaire pour tout $n \ge 1$.\\
        Par le principe de récurrence, pour tout $n \in \N^*$ :
        \[
        OA_n = \sqrt{n}
        \]
    \end{itemize}

    \item \textbf{Synthèse géométrique et figure d'ensemble :} \bareme{1,0 pt}
\end{enumerate}
\end{correctionbox}

\begin{prolongementbox}[La Spirale de Théodore dans l'Histoire et le dialogue du \emph{Théétète}]
La construction de cette figure remonte au philosophe et mathématicien présocratique \textbf{Théodore de Cyrène} (V\ieme{} siècle av. J.-C.), maître de mathématiques de Platon. Dans le célèbre dialogue platonicien du \emph{Théétète} (147d--148b), Théodore expose la démonstration de l'incommensurabilité (l'irrationalité) des racines carrées d'entiers qui ne sont pas des carrés parfaits :
\begin{quote}
\og Théodore nous avait tracé des figures pour nous faire voir que les côtés des carrés de 3 pieds et de 5 pieds ne sont pas commensurables en longueur avec celui d'un pied, et il prit ainsi chaque cas particulier jusqu'à celui de 17 pieds, où il s'arrêta. \fg
\end{quote}

\textbf{Pourquoi Théodore s'est-il arrêté précisément à $\sqrt{17}$ ?}\\
La réponse tient à la géométrie même de la spirale : à chaque étape, l'angle au sommet $O$ du triangle rectangle vaut :
\[
\alpha_k = \widehat{A_k O A_{k+1}} = \arctan\left(\frac{A_k A_{k+1}}{OA_k}\right) = \arctan\left(\frac{1}{\sqrt{k}}\right)
\]
L'angle polaire total cumulé après $n$ triangles vaut donc $\Theta(n) = \sum_{k=1}^n \arctan\left(\frac{1}{\sqrt{k}}\right)$.\\
Or :
\begin{align*}
\Theta(16) &= \sum_{k=1}^{16} \arctan\left(\frac{1}{\sqrt{k}}\right) \approx 351{,}15^\circ < 360^\circ = 2\pi\text{ rad} \\
\Theta(17) &= \Theta(16) + \arctan\left(\frac{1}{\sqrt{17}}\right) \approx 351{,}15^\circ + 13{,}63^\circ = 364{,}78^\circ > 360^\circ
\end{align*}
C'est donc au 17\ieme{} triangle ($OA_{16}A_{17}$) que la spirale effectue son premier tour complet sans jamais se croiser. Ajouter un 18\ieme{} triangle ferait chevaucher le segment initial $[OA_1]$ !

\begin{center}
\begin{tikzpicture}[scale=1.1, >=latex, line cap=round, line join=round]
  \coordinate (O) at (0,0);
  \fill[black] (O) circle (1.8pt);
  \node[below left=2pt, font=\small] at (O) {$O$};

  % Coordonnées de la spirale complète (s = 1.05 cm)
  \coordinate (B1) at (1.050, 0.000);
  \coordinate (B2) at (1.050, 1.050);
  \coordinate (B3) at (0.308, 1.792);
  \coordinate (B4) at (-0.727, 1.970);
  \coordinate (B5) at (-1.712, 1.606);
  \coordinate (B6) at (-2.431, 0.841);
  \coordinate (B7) at (-2.774, -0.152);
  \coordinate (B8) at (-2.717, -1.200);
  \coordinate (B9) at (-2.292, -2.161);
  \coordinate (B10) at (-1.572, -2.925);
  \coordinate (B11) at (-0.647, -3.422);
  \coordinate (B12) at (0.385, -3.617);
  \coordinate (B13) at (1.429, -3.506);
  \coordinate (B14) at (2.401, -3.110);
  \coordinate (B15) at (3.232, -2.468);
  \coordinate (B16) at (3.869, -1.633);
  \coordinate (B17) at (4.278, -0.666);

  % Remplissage en dégradé subtil
  \foreach \i/\j/\c in {1/2/5, 2/3/9, 3/4/13, 4/5/17, 5/6/21, 6/7/25, 7/8/29, 8/9/33, 9/10/37, 10/11/41, 11/12/45, 12/13/49, 13/14/53, 14/15/57, 15/16/61, 16/17/65} {
    \fill[royalpurple!\c!white] (O) -- (B\i) -- (B\j) -- cycle;
    \draw[thin, black!35] (O) -- (B\j);
    \draw[thick, secondary] (B\i) -- (B\j);
  }
  \draw[thick, primary] (O) -- (B1);

  % Étiquettes des hypoténuses
  \node[below, font=\tiny] at ($(O)!0.5!(B1)$) {$1$};
  \node[above left=-2pt, font=\tiny] at ($(O)!0.5!(B2)$) {$\sqrt{2}$};
  \node[above=1pt, font=\tiny] at ($(O)!0.5!(B3)$) {$\sqrt{3}$};
  \node[left=1pt, font=\tiny] at ($(O)!0.5!(B4)$) {$\sqrt{4}$};
  \node[left=1pt, font=\tiny] at ($(O)!0.5!(B5)$) {$\sqrt{5}$};
  \node[left=1pt, font=\tiny] at ($(O)!0.5!(B6)$) {$\sqrt{6}$};
  \node[left=1pt, font=\tiny] at ($(O)!0.5!(B7)$) {$\sqrt{7}$};
  \node[below left=1pt, font=\tiny] at ($(O)!0.5!(B8)$) {$\sqrt{8}$};
  \node[below left=1pt, font=\tiny] at ($(O)!0.5!(B9)$) {$\sqrt{9}$};
  \node[below=1pt, font=\tiny] at ($(O)!0.5!(B10)$) {$\sqrt{10}$};
  \node[below=1pt, font=\tiny] at ($(O)!0.5!(B11)$) {$\sqrt{11}$};
  \node[below=1pt, font=\tiny] at ($(O)!0.5!(B12)$) {$\sqrt{12}$};
  \node[below right=1pt, font=\tiny] at ($(O)!0.5!(B13)$) {$\sqrt{13}$};
  \node[below right=1pt, font=\tiny] at ($(O)!0.5!(B14)$) {$\sqrt{14}$};
  \node[right=1pt, font=\tiny] at ($(O)!0.5!(B15)$) {$\sqrt{15}$};
  \node[right=1pt, font=\tiny] at ($(O)!0.5!(B16)$) {$\sqrt{16}$};
  \node[right=1pt, font=\tiny] at ($(O)!0.5!(B17)$) {$\sqrt{17}$};

  % Sommets
  \foreach \k in {1,2,3,4,5,6,7,8,9,10,11,12,13,14,15,16,17} {
    \fill[black] (B\k) circle (1.2pt);
  }
  \node[right, font=\scriptsize] at (B1) {$A_1$};
  \node[above right, font=\scriptsize] at (B2) {$A_2$};
  \node[above, font=\scriptsize] at (B4) {$A_4$};
  \node[left, font=\scriptsize] at (B7) {$A_7$};
  \node[below left, font=\scriptsize] at (B10) {$A_{10}$};
  \node[below, font=\scriptsize] at (B12) {$A_{12}$};
  \node[below right, font=\scriptsize] at (B14) {$A_{14}$};
  \node[right, font=\scriptsize] at (B17) {$A_{17}$};

  % Arc de cercle indiquant la fermeture angulaire à 351°
  \draw[->, thick, primary, dashed] (0.8, 0) arc (0:351.15:0.8);
  \node[above right, font=\tiny, primary] at (0.85, 0.2) {$\Theta(16) \approx 351{,}15^\circ$};
\end{tikzpicture}
\end{center}
\end{prolongementbox}

\vspace{0.3cm}

% ----------------------------------------------------
% CORRIGÉ EXERCICE 3
% ----------------------------------------------------
\subsection*{Correction détaillée de l'Exercice 3}

\begin{correctionbox}[Corrigé de l'Exercice 3 \bareme{Total : 4 points}]
La suite $(u_n)_{n \in \N}$ est définie par $u_0 = 0$ et pour tout $n \in \N$, $u_{n+1} = u_n + 3n(n+1) + 1$.

\begin{enumerate}[label=\textbf{\arabic*.}]
    \item \textbf{Calcul des premiers termes $u_1$, $u_2$ et $u_3$ :} \bareme{1,0 pt}
    \begin{itemize}[leftmargin=*, label=$\bullet$]
        \item Pour $n = 0$ :
        \[
        u_1 = u_0 + 3 \times 0 \times (0+1) + 1 = 0 + 0 + 1 = 1 \bareme{0,25 pt}
        \]
        \item Pour $n = 1$ :
        \[
        u_2 = u_1 + 3 \times 1 \times (1+1) + 1 = 1 + 3 \times 2 + 1 = 1 + 6 + 1 = 8 \bareme{0,25 pt}
        \]
        \item Pour $n = 2$ :
        \[
        u_3 = u_2 + 3 \times 2 \times (2+1) + 1 = 8 + 3 \times 6 + 1 = 8 + 18 + 1 = 27 \bareme{0,5 pt}
        \]
    \end{itemize}

    \item \textbf{Conjecture de l'expression de $u_n$ en fonction de $n$ :} \bareme{0,5 pt}\\
    Dressons le tableau des premiers termes :
    \begin{center}
    \begin{tabular}{|c|c|c|c|c|}
    \hline
    $n$ & 0 & 1 & 2 & 3 \\
    \hline
    $u_n$ & $0$ & $1$ & $8$ & $27$ \\
    \hline
    $n^3$ & $0^3 = 0$ & $1^3 = 1$ & $2^3 = 8$ & $3^3 = 27$ \\
    \hline
    \end{tabular}
    \end{center}
    On observe que pour $n \in \{0, 1, 2, 3\}$, on a $u_n = n^3$. On conjecture donc que pour tout $n \in \N$ :
    \[
    u_n = n^3
    \]

    \item \textbf{Démonstration de la conjecture par récurrence :} \bareme{2,5 pts}

    Pour tout $n \in \N$, notons le prédicat :
    \[
    \mathcal{P}(n) : \quad \ll u_n = n^3 \gg
    \]

    \begin{itemize}[leftmargin=*, label=$\bullet$]
        \item \textbf{Initialisation ($n = 0$) :} \pts{0,5 pt}\\
        D'une part, $u_0 = 0$.\\
        D'autre part, $0^3 = 0$.\\
        L'égalité $u_0 = 0^3$ est donc satisfaite. La propriété $\mathcal{P}(0)$ est vraie.

        \item \textbf{Hérédité :} \pts{1,5 pt}\\
        Soit $k \in \N$ un entier naturel fixé.\\
        Supposons que $\mathcal{P}(k)$ soit vraie, à savoir :
        \[
        u_k = k^3 \quad \text{(\textbf{Hypothèse de Récurrence})}
        \]
        Montrons que $\mathcal{P}(k+1)$ est vraie, à savoir :
        \[
        u_{k+1} = (k+1)^3
        \]
        Par la relation de récurrence liant les termes de la suite :
        \[
        u_{k+1} = u_k + 3k(k+1) + 1
        \]
        En substituant $u_k$ par sa valeur $k^3$ donnée par l'hypothèse de récurrence :
        \[
        u_{k+1} = k^3 + 3k(k+1) + 1
        \]
        Développons l'expression algébrique :
        \[
        u_{k+1} = k^3 + 3k^2 + 3k + 1
        \]
        Or, développons le cube $(k+1)^3$ à l'aide de la formule du binôme de Newton (ou par distributivité usuelle) :
        \[
        (k+1)^3 = (k+1)(k+1)^2 = (k+1)(k^2 + 2k + 1) = k^3 + 2k^2 + k + k^2 + 2k + 1 = k^3 + 3k^2 + 3k + 1
        \]
        On constate l'égalité parfaite des deux expressions :
        \[
        u_{k+1} = (k+1)^3
        \]
        La proposition $\mathcal{P}(k+1)$ est ainsi établie. L'hérédité est prouvée pour tout entier $k \in \N$.

        \item \textbf{Conclusion :} \pts{0,5 pt}\\
        La propriété est initialisée au rang $0$ et est héréditaire.\\
        D'après le principe de raisonnement par récurrence, pour tout $n \in \N$ :
        \[
        u_n = n^3
        \]
    \end{itemize}
\end{enumerate}
\end{correctionbox}

\begin{vigilancebox}
\textbf{Conseil de calcul :} La formule du binôme de Newton $(a+b)^3 = a^3 + 3a^2b + 3ab^2 + b^3$ doit être maîtrisée sans hésitation en classe préparatoire. Savoir reconnaître immédiatement le développement $k^3 + 3k^2 + 3k + 1 = (k+1)^3$ fait gagner un temps précieux et évite les erreurs de signe.
\end{vigilancebox}

\begin{prolongementbox}[Formule de Faulhaber et calcul de la somme des carrés]
L'exercice 3 fournit une méthode directe pour calculer la somme des carrés $\sum_{k=1}^n k^2$ sans connaître la formule à l'avance !\\
En effet, la relation $u_{k+1} - u_k = 3k(k+1) + 1 = 3k^2 + 3k + 1$ s'écrit :
\[
(k+1)^3 - k^3 = 3k^2 + 3k + 1
\]
En sommant cette identité pour $k = 1$ jusqu'à $n$ :
\[
\sum_{k=1}^n \left( (k+1)^3 - k^3 \right) = 3 \sum_{k=1}^n k^2 + 3 \sum_{k=1}^n k + \sum_{k=1}^n 1
\]
Par télescopage, le membre de gauche vaut $(n+1)^3 - 1^3$. Sachant que $\sum_{k=1}^n k = \dfrac{n(n+1)}{2}$ et $\sum_{k=1}^n 1 = n$ :
\[
(n+1)^3 - 1 = 3 \sum_{k=1}^n k^2 + 3 \frac{n(n+1)}{2} + n
\]
En isolant $\sum_{k=1}^n k^2$, on retrouve la célèbre \textbf{formule de Faulhaber} pour les carrés :
\[
\sum_{k=1}^n k^2 = \frac{n(n+1)(2n+1)}{6}
\]
Ce principe se généralise à toutes les sommes de puissances $\sum_{k=1}^n k^p$ via les \textbf{polynômes de Bernoulli}.
\end{prolongementbox}

\vspace{0.3cm}

% ----------------------------------------------------
% CORRIGÉ EXERCICE 4
% ----------------------------------------------------
\subsection*{Correction détaillée de l'Exercice 4}

\begin{correctionbox}[Corrigé de l'Exercice 4 \bareme{Total : 4 points}]
\textbf{Énoncé :} Pour $n \in \N^*$, on souhaite démontrer par récurrence que :
\[
\sum_{i=1}^n (-1)^i \, i = \frac{(-1)^n (2n+1) - 1}{4}
\]

Pour tout entier naturel non nul $n \in \N^*$, posons le prédicat :
\[
\mathcal{P}(n) : \quad \ll \sum_{i=1}^n (-1)^i \, i = \frac{(-1)^n (2n+1) - 1}{4} \gg
\]

\begin{itemize}[leftmargin=*, label=$\bullet$]
    \item \textbf{Initialisation ($n = 1$) :} \bareme{1,0 pt}
    \begin{itemize}[leftmargin=*, label=$-$]
        \item Calculons la valeur du membre de gauche pour $n = 1$ :
        \[
        \sum_{i=1}^1 (-1)^i \, i = (-1)^1 \times 1 = -1 \times 1 = -1
        \]
        \item Calculons la valeur du membre de droite pour $n = 1$ :
        \[
        \frac{(-1)^1 \big(2(1) + 1\big) - 1}{4} = \frac{-1 \times 3 - 1}{4} = \frac{-3 - 1}{4} = \frac{-4}{4} = -1
        \]
    \end{itemize}
    Les deux membres coïncident rigoureusement : $-1 = -1$.\\
    Le prédicat $\mathcal{P}(1)$ est donc vérifié.

    \item \textbf{Hérédité :} \bareme{2,5 pts}\\
    Soit $k \in \N^*$ un entier quelconque fixé.\\
    Supposons que la propriété $\mathcal{P}(k)$ soit vraie, c'est-à-dire :
    \[
    \sum_{i=1}^k (-1)^i \, i = \frac{(-1)^k (2k+1) - 1}{4} \quad \text{(\textbf{Hypothèse de Récurrence})}
    \]
    Montrons que $\mathcal{P}(k+1)$ est vraie, c'est-à-dire que :
    \[
    \sum_{i=1}^{k+1} (-1)^i \, i = \frac{(-1)^{k+1} \big(2(k+1)+1\big) - 1}{4} = \frac{(-1)^{k+1} (2k+3) - 1}{4}
    \]
    
    Décomposons la somme au rang $k+1$ en isolant son tout dernier terme :
    \[
    \sum_{i=1}^{k+1} (-1)^i \, i = \left(\sum_{i=1}^k (-1)^i \, i\right) + (-1)^{k+1} (k+1)
    \]
    Par application directe de l'hypothèse de récurrence :
    \[
    \sum_{i=1}^{k+1} (-1)^i \, i = \frac{(-1)^k (2k+1) - 1}{4} + (-1)^{k+1} (k+1)
    \]
    Réduisons au même dénominateur $4$ :
    \[
    \sum_{i=1}^{k+1} (-1)^i \, i = \frac{(-1)^k (2k+1) - 1 + 4 \cdot (-1)^{k+1} (k+1)}{4}
    \]
    Utilisons la relation algébrique fondamentale sur les signes alternés :
    \[
    (-1)^k = - (-1)^{k+1} \quad \left(\text{en effet, } (-1)^{k+1} = (-1) \times (-1)^k \iff (-1)^k = -(-1)^{k+1}\right)
    \]
    Substituons cette égalité dans le numérateur :
    \begin{align*}
    \text{Numérateur} &= \Big[ -(-1)^{k+1} \Big] (2k+1) - 1 + 4 \cdot (-1)^{k+1} (k+1) \\
                      &= (-1)^{k+1} \Big[ -(2k+1) + 4(k+1) \Big] - 1 \\
                      &= (-1)^{k+1} \Big[ -2k - 1 + 4k + 4 \Big] - 1 \\
                      &= (-1)^{k+1} (2k + 3) - 1
    \end{align*}
    Or, remarquons que $2k + 3 = 2(k+1) + 1$.\\
    Par conséquent :
    \[
    \sum_{i=1}^{k+1} (-1)^i \, i = \frac{(-1)^{k+1} \big(2(k+1)+1\big) - 1}{4}
    \]
    La proposition $\mathcal{P}(k+1)$ est ainsi établie. L'hérédité est démontrée pour tout $k \in \N^*$.

    \item \textbf{Conclusion :} \bareme{0,5 pt}\\
    La propriété $\mathcal{P}(1)$ est vraie et la transmission d'hérédité est assurée pour tout entier $n \ge 1$.\\
    En vertu du principe de récurrence, pour tout $n \in \N^*$ :
    \[
    \sum_{i=1}^n (-1)^i \, i = \frac{(-1)^n (2n+1) - 1}{4}
    \]
\end{itemize}
\end{correctionbox}

\begin{vigilancebox}
\textbf{Technique de factorisation de $(-1)^n$ :} L'erreur majeure des élèves consiste à développer sans factoriser les puissances de $(-1)$, conduisant à des blocages algébriques insolubles. Il faut systématiquement convertir $(-1)^k$ en $-(-1)^{k+1}$ (ou réciproquement factoriser par $(-1)^k$) pour mettre en facteur commun la puissance au rang cible.
\end{vigilancebox}

\begin{prolongementbox}[Démonstration alternative sans récurrence : Parité \& Analyse]
\textbf{1. Approche arithmétique par regroupement selon la parité :}
Écrivons la somme $S_n = \sum_{i=1}^n (-1)^i i = -1 + 2 - 3 + 4 - 5 + \dots + (-1)^n n$.
\begin{itemize}[leftmargin=*, label=$\bullet$]
    \item Si $n = 2m$ est pair, groupons les termes deux par deux :
    \[
    S_{2m} = (-1 + 2) + (-3 + 4) + \dots + (-(2m-1) + 2m) = \underbrace{1 + 1 + \dots + 1}_{m \text{ termes}} = m = \frac{n}{2}
    \]
    La formule donne : $\dfrac{(-1)^{2m}(2(2m)+1)-1}{4} = \dfrac{1 \cdot (4m+1)-1}{4} = \dfrac{4m}{4} = m$. La formule est vérifiée !
    \item Si $n = 2m+1$ est impair :
    \[
    S_{2m+1} = S_{2m} + (-1)^{2m+1}(2m+1) = m - (2m+1) = -m - 1 = -\frac{n+1}{2}
    \]
    La formule donne : $\dfrac{(-1)^{2m+1}(2(2m+1)+1)-1}{4} = \dfrac{-(4m+3)-1}{4} = \dfrac{-4m-4}{4} = -(m+1)$. Concordance parfaite !
\end{itemize}

\textbf{2. Approche analytique par dérivation de série géométrique (Outil de Sup) :}\\
Pour $x \neq 1$, considérons la somme géométrique $\sum_{i=0}^n x^i = \dfrac{1 - x^{n+1}}{1 - x}$.\\
En dérivant terme à terme par rapport à $x$ :
\[
\sum_{i=1}^n i x^{i-1} = \frac{-(n+1)x^n(1-x) - (1-x^{n+1})(-1)}{(1-x)^2} = \frac{-(n+1)x^n + (n+1)x^{n+1} + 1 - x^{n+1}}{(1-x)^2}
\]
En multipliant par $x$ :
\[
\sum_{i=1}^n i x^i = \frac{-(n+1)x^{n+1} + n x^{n+2} + x}{(1-x)^2}
\]
En évaluant cette fraction rationnelle en $x = -1$ (le dénominateur valant $(1 - (-1))^2 = 4$) :
\[
\sum_{i=1}^n (-1)^i i = \frac{-(n+1)(-1)^{n+1} + n(-1)^{n+2} - 1}{4} = \frac{(-1)^n \Big[ (n+1) + n \Big] - 1}{4} = \frac{(-1)^n (2n+1) - 1}{4}
\]
On retrouve la formule en deux lignes de calcul différentiel formel !
\end{prolongementbox}

\vspace{0.3cm}

% ----------------------------------------------------
% CORRIGÉ EXERCICE 5
% ----------------------------------------------------
\subsection*{Correction détaillée de l'Exercice 5}

\begin{correctionbox}[Corrigé de l'Exercice 5 \bareme{Total : 4 points}]
\begin{enumerate}[label=\textbf{\arabic*.}]
    \item \textbf{Démontrons par récurrence que pour tout $n \in \N$, il existe $b_n \in \N$ tel que $10^n = 3b_n + 1$.} \bareme{2,0 pts}

    Pour tout entier naturel $n \in \N$, posons la propriété :
    \[
    \mathcal{P}(n) : \quad \ll \exists b_n \in \N, \quad 10^n = 3b_n + 1 \gg
    \]

    \begin{itemize}[leftmargin=*, label=$\bullet$]
        \item \textbf{Initialisation ($n = 0$) :} \pts{0,5 pt}\\
        Pour $n = 0$, $10^0 = 1$.\\
        Or, on peut écrire :
        \[
        1 = 3 \times 0 + 1
        \]
        Puisque $0 \in \N$, en posant $b_0 = 0$, on a bien $b_0 \in \N$ et $10^0 = 3b_0 + 1$.\\
        La proposition $\mathcal{P}(0)$ est donc vraie.

        \item \textbf{Hérédité :} \pts{1,25 pt}\\
        Soit $k \in \N$ un entier naturel fixé.\\
        Supposons que $\mathcal{P}(k)$ soit vraie, c'est-à-dire qu'il existe un entier $b_k \in \N$ vérifiant :
        \[
        10^k = 3b_k + 1 \quad \text{(\textbf{Hypothèse de Récurrence})}
        \]
        Montrons que $\mathcal{P}(k+1)$ est vraie, c'est-à-dire qu'il existe un entier $b_{k+1} \in \N$ tel que :
        \[
        10^{k+1} = 3b_{k+1} + 1
        \]
        Exprimons $10^{k+1}$ en fonction de $10^k$ :
        \[
        10^{k+1} = 10 \times 10^k
        \]
        Remplaçons $10^k$ par l'expression issue de l'hypothèse de récurrence :
        \begin{align*}
        10^{k+1} &= 10(3b_k + 1) \\
                 &= 30 b_k + 10
        \end{align*}
        Décomposons astucieusement l'entier $10$ sous la forme $9 + 1 = 3 \times 3 + 1$ afin de faire émerger le reste $1$ :
        \begin{align*}
        10^{k+1} &= 30 b_k + 9 + 1 \\
                 &= 3(10 b_k + 3) + 1
        \end{align*}
        Posons :
        \[
        b_{k+1} = 10 b_k + 3
        \]
        Puisque $b_k \in \N$, par produit et somme d'entiers naturels, on a de manière certaine $b_{k+1} = 10b_k + 3 \in \N$.\\
        Nous avons ainsi prouvé l'existence d'un entier naturel $b_{k+1} \in \N$ tel que $10^{k+1} = 3b_{k+1} + 1$.\\
        La propriété $\mathcal{P}(k+1)$ est donc vraie. L'hérédité est validée.

        \item \textbf{Conclusion :} \pts{0,25 pt}\\
        Par le principe de récurrence, pour tout $n \in \N$, il existe un entier $b_n \in \N$ tel que :
        \[
        10^n = 3b_n + 1
        \]
        \emph{Remarque constructive :} On observe que $b_0 = 0$, $b_1 = 3$, $b_2 = 33$, $b_3 = 333$ et d'une manière générale, $b_n = \underbrace{33\dots3}_{n \text{ fois}} = \dfrac{10^n - 1}{3}$.
    \end{itemize}

    \item \textbf{Démontrons l'équivalence de divisibilité par $3$ et interprétation :} \bareme{2,0 pts}

    Soit $p \in \N$ et soient $a_0, a_1, \dots, a_p$ des entiers naturels (représentant les chiffres de l'écriture en base $10$).\\
    Considérons la somme $\dsp N = \sum_{k=0}^p a_k 10^k$.\\
    D'après le résultat de la question 1, pour chaque indice $k \in \{0, \dots, p\}$, il existe $b_k \in \N$ tel que $10^k = 3b_k + 1$.\\
    Substituons cette identité dans la somme $N$ :
    \[
    \sum_{k=0}^p a_k 10^k = \sum_{k=0}^p a_k (3b_k + 1) = \sum_{k=0}^p \left( 3 a_k b_k + a_k \right)
    \]
    Par linéarité de la somme finie :
    \[
    \sum_{k=0}^p a_k 10^k = 3 \left( \sum_{k=0}^p a_k b_k \right) + \sum_{k=0}^p a_k
    \]
    Posons $M = \sum_{k=0}^p a_k b_k$. Puisque chaque $a_k \in \N$ et $b_k \in \N$, $M$ est un entier naturel ($M \in \N$).\\
    L'égalité s'écrit de façon compacte :
    \[
    \sum_{k=0}^p a_k 10^k - \sum_{k=0}^p a_k = 3M
    \]
    Ce qui prouve que la différence entre le nombre $\sum_{k=0}^p a_k 10^k$ et la somme de ses coefficients $\sum_{k=0}^p a_k$ est un \textbf{multiple de $3$}. \pts{1,0 pt}

    Établissons à présent la double implication :
    \begin{itemize}[leftmargin=*, label=$\bullet$]
        \item \textbf{Sens direct ($\implies$) :} Supposons que $3$ divise $\sum_{k=0}^p a_k 10^k$. Alors il existe un entier $q \in \Z$ tel que $\sum_{k=0}^p a_k 10^k = 3q$. Dès lors :
        \[
        \sum_{k=0}^p a_k = \sum_{k=0}^p a_k 10^k - 3M = 3q - 3M = 3(q - M)
        \]
        Puisque $(q - M) \in \Z$, $3$ divise $\sum_{k=0}^p a_k$.
        \item \textbf{Sens réciproque ($\impliedby$) :} Réciproquement, supposons que $3$ divise $\sum_{k=0}^p a_k$. Alors il existe $q' \in \Z$ tel que $\sum_{k=0}^p a_k = 3q'$. Dès lors :
        \[
        \sum_{k=0}^p a_k 10^k = 3M + \sum_{k=0}^p a_k = 3M + 3q' = 3(M + q')
        \]
        Puisque $(M + q') \in \Z$, $3$ divise $\sum_{k=0}^p a_k 10^k$.
    \end{itemize}
    On a donc démontré rigoureusement l'équivalence :
    \[
    \left( 3 \text{ divise } \sum_{k=0}^p a_k 10^k \right) \iff \left( 3 \text{ divise } \sum_{k=0}^p a_k \right) \bareme{0,5 pt}
    \]

    \textbf{Interprétation arithmétique :} \bareme{0,5 pt}\\
    Tout entier naturel $N$ non nul s'écrit de manière unique en base décimale sous la forme :
    \[
    N = a_p 10^p + a_{p-1} 10^{p-1} + \dots + a_1 10 + a_0 = \sum_{k=0}^p a_k 10^k
    \]
    où les $a_k \in \{0, 1, 2, \dots, 9\}$ sont les chiffres de son écriture usuelle en base 10.\\
    Le terme $\sum_{k=0}^p a_k$ est précisément la \textbf{somme des chiffres décimaux} du nombre $N$.\\
    L'équivalence obtenue démontre donc le résultat fondamental de l'arithmétique élémentaire :
    \begin{center}
    \textbf{\og Un nombre entier naturel est divisible par $3$ si et seulement si la somme de ses chiffres est divisible par $3$. \fg}
    \end{center}
\end{enumerate}
\end{correctionbox}

\begin{prolongementbox}[Arithmétique modulaire de Gauss et généralisation des critères]
Dans le langage moderne des \textbf{congruences} forgé par Carl Friedrich Gauss dans ses \emph{Disquisitiones Arithmeticae} (1801) :
\[
10 = 3 \times 3 + 1 \implies 10 \equiv 1 \pmod 3
\]
Par compatibilité élémentaire de la congruence avec les puissances entières :
\[
\forall k \in \N, \quad 10^k \equiv 1^k \equiv 1 \pmod 3
\]
Par conséquent :
\[
\sum_{k=0}^p a_k 10^k \equiv \sum_{k=0}^p a_k (1) \equiv \sum_{k=0}^p a_k \pmod 3
\]
L'équivalence de divisibilité est donc un corollaire direct et immédiat de la structure d'anneau du quotient $\Z/3\Z$.

\textbf{Généralisations universelles en base 10 :}
\begin{enumerate}[leftmargin=*, label=$\bullet$]
    \item \textbf{Critère par 9 :} Comme $10 = 9 \times 1 + 1$, $10 \equiv 1 \pmod 9 \implies 10^k \equiv 1 \pmod 9$. Un nombre est divisible par $9$ si et seulement si la somme de ses chiffres est divisible par $9$.
    \item \textbf{Critère par 11 :} Comme $10 = 11 - 1$, on a $10 \equiv -1 \pmod{11}$. Par conséquent :
    \[
    10^k \equiv (-1)^k \pmod{11} \implies \sum_{k=0}^p a_k 10^k \equiv \sum_{k=0}^p (-1)^k a_k \pmod{11}
    \]
    Un nombre est divisible par $11$ si et seulement si la \textbf{somme alternée de ses chiffres} ($a_0 - a_1 + a_2 - a_3 + \dots$) est un multiple de $11$.
    \item \textbf{Critère par 7 et 13 (Le Ruban de Pascal) :} Comme $1000 = 10^3 = 7 \times 143 - 1 = 13 \times 77 - 1$, on a $10^3 \equiv -1 \pmod 7$ et $10^3 \equiv -1 \pmod{13}$. Il suffit donc de faire la somme alternée des \textbf{tranches de trois chiffres} de droite à gauche !
\end{enumerate}
\end{prolongementbox}

\vspace{0.3cm}

% ----------------------------------------------------
% CORRIGÉ DU BONUS
% ----------------------------------------------------
\subsection*{Correction détaillée du Bonus}

\begin{correctionbox}[Corrigé du Bonus (+2 points) \bareme{Total : +2 points}]
\textbf{Objectif :} Démontrer la relation $u_n = n^3$ pour tout $n \in \N$ de façon directe, sans récurrence.

\medskip
Reprenons la relation de définition de la suite de l'exercice 3 :
\[
\begin{cases}
u_0 = 0 \\
u_{k+1} - u_k = 3k(k+1) + 1 \quad (k \in \N)
\end{cases}
\]
Développons le membre de droite :
\[
3k(k+1) + 1 = 3k^2 + 3k + 1
\]
Or, considérons l'identité algébrique remarquable du cube de deux entiers consécutifs :
\[
(k+1)^3 - k^3 = \big(k^3 + 3k^2 + 3k + 1\big) - k^3 = 3k^2 + 3k + 1
\]
On observe donc une identité remarquable exacte terme à terme :
\[
u_{k+1} - u_k = (k+1)^3 - k^3 \bareme{0,75 pt}
\]

Soit $n \in \N^*$ un entier naturel fixé. Sommons cette égalité pour l'indice muet $k$ variant de $0$ jusqu'à $n-1$ :
\[
\sum_{k=0}^{n-1} (u_{k+1} - u_k) = \sum_{k=0}^{n-1} \left[ (k+1)^3 - k^3 \right] \bareme{0,5 pt}
\]
Ces deux sommes sont des \textbf{sommes télescopiques} où tous les termes intermédiaires s'annulent deux à deux :
\begin{itemize}[leftmargin=*, label=$\bullet$]
    \item Pour le membre de gauche :
    \[
    \sum_{k=0}^{n-1} (u_{k+1} - u_k) = (u_1 - u_0) + (u_2 - u_1) + (u_3 - u_2) + \dots + (u_n - u_{n-1}) = u_n - u_0
    \]
    \item Pour le membre de droite :
    \[
    \sum_{k=0}^{n-1} \left[ (k+1)^3 - k^3 \right] = (1^3 - 0^3) + (2^3 - 1^3) + (3^3 - 2^3) + \dots + (n^3 - (n-1)^3) = n^3 - 0^3 = n^3
    \]
\end{itemize}
En égalant les deux bilans télescopiques :
\[
u_n - u_0 = n^3
\]
Comme la condition initiale est $u_0 = 0$, on obtient directement :
\[
u_n = n^3 \bareme{0,75 pt}
\]
Enfin, pour $n = 0$, $u_0 = 0$ et $0^3 = 0$, l'égalité est également vérifiée au rang 0.

L'égalité $u_n = n^3$ est donc établie directement pour tout $n \in \N$ sans aucune induction.
\end{correctionbox}

\vspace{0.4cm}

% ----------------------------------------------------
% FICHE DE SYNTHÈSE MÉTHODOLOGIQUE MPSI
% ----------------------------------------------------
\begin{tcolorbox}[enhanced, colback=white, colframe=primary, arc=2mm, boxrule=1pt, title={\textsc{Fiche Méthode MPSI : Les Trois Piliers du Raisonnement par Récurrence}}]
\small
Le raisonnement par récurrence repose sur la propriété du \textbf{bon ordre} de l'ensemble $\N$ (tout sous-ensemble non vide de $\N$ admet un plus petit élément), axiome fondamental formalisé par Giuseppe Peano (1889).

Pour obtenir la note maximale aux concours et en prépa, toute copie doit respecter scrupuleusement le canevas en trois phases :

\begin{enumerate}[leftmargin=*]
    \item \textbf{Déclaration claire du prédicat :}
    Ne jamais commencer par écrire des calculs dans le vide. Toujours poser explicitement :
    \begin{center}
    \og Pour tout $n \ge n_0$, notons $\mathcal{P}(n)$ la proposition : $\ll \dots \gg$. \fg
    \end{center}

    \item \textbf{L'Initialisation (Phase d'ancrage) :}
    Évaluer séparément le membre de gauche et le membre de droite pour le plus petit rang $n_0$ (généralement $n_0 = 0$ ou $n_0 = 1$). Conclure par \og Donc $\mathcal{P}(n_0)$ est vraie. \fg

    \item \textbf{L'Hérédité (Phase de transmission) :}
    La rédaction doit être sans faille :
    \begin{center}
    \og \emph{Soit $k \ge n_0$ un entier fixé. Supposons $\mathcal{P}(k)$ vraie. Montrons que $\mathcal{P}(k+1)$ est vraie.} \fg
    \end{center}
    Dans le corps de la preuve, expliciter clairement l'étape où l'hypothèse de récurrence est injectée (mentionner \textbf{[D'après l'H.R.]}).

    \item \textbf{La Conclusion logique :}
    Rappeler que les deux conditions de l'axiome sont réunies :
    \begin{center}
    \og \emph{La proposition est initialisée en $n_0$ et est héréditaire. Par le principe de raisonnement par récurrence, $\mathcal{P}(n)$ est vraie pour tout entier $n \ge n_0$.} \fg
    \end{center}
\end{enumerate}
\end{tcolorbox}

\end{document}
